% Layout fixture: a drawing set with LaTeX's picture environment. The slanted side of the triangle
% is a row of glyphs of the font line10 (each a short piece of the line), not a path: it is one
% figure, not a paragraph of quote marks (F-550). The paragraphs around it are plain text.
% Written for ReadableSTEM's tests; no third-party text.
\documentclass[11pt]{article}
\usepackage[OT1]{fontenc}
\renewcommand{\rmdefault}{cmr}
\usepackage[margin=1in]{geometry}
% The layout reads TeX conventions only from a PDF that says TeX made it (/Creator); Tectonic's
% own creator string does not, so the fixture sets one as a pdfLaTeX file has it.
\usepackage[pdfcreator={LaTeX with hyperref}]{hyperref}
\pagestyle{empty}
\setlength{\unitlength}{1mm}

\begin{document}

The ramp in the next figure rises 3.5 m over a horizontal run of 6 m. Use it to find the length
of the slope, and say what angle the slope makes with the ground.

\begin{center}
\begin{picture}(80,45)
 \put(10,5){\line(1,0){60}} \put(70,5){\line(0,1){35}} \put(10,5){\line(5,3){60}}
 \put(38,0){$a = 6$} \put(72,20){$b = 3.5$} \put(30,25){$c$}
 \put(15,7){$\theta$} \put(66,7){\framebox(4,4){}}
 \put(5,2){$A$} \put(71,2){$B$} \put(71,41){$C$}
\end{picture}

Figure 1: A right triangle with legs $a$ and $b$.
\end{center}

The slope is the longest side of the triangle, and the ground is the shortest horizontal one.

\end{document}
